\documentclass[10pt]{article}
\usepackage{icassp2027_paperkit}
\usepackage{amsmath,amssymb,booktabs,graphicx,microtype}
\usepackage{url}
\urlstyle{same}
\usepackage{placeins}
\usepackage{cuted}
\usepackage{balance}
\usepackage[T1]{fontenc}
\usepackage{textcomp}
\setlength{\textfloatsep}{6pt plus 2pt minus 2pt}
\setlength{\floatsep}{5pt plus 2pt minus 2pt}
\setlength{\intextsep}{5pt plus 2pt minus 2pt}
\setlength{\abovecaptionskip}{2pt}
\setlength{\belowcaptionskip}{0pt}

\newcommand{\dropv}{d}
\newcommand{\foils}{\mathcal{F}}
\newcommand{\regions}{\mathcal{R}}

\title{Beyond Mean Foils: Auditing Worst-Foil Specificity in Frozen CLIP Region Explanations}
% ICASSP 2027 single-anonymous author block.
% Keep this order identical to the online submission system.
\name{Kaixin Liu$^{1,*}$, Zhipeng Ye$^{1,*,\dagger}$, Feng Jiang$^{1}$, Qiufeng Wang$^{2}$}
\address{$^{1}$Taizhou Institute of Science and Technology, Nanjing University of Science and Technology,\\
Taizhou 225300, Jiangsu, China\\
$^{2}$Department of Intelligence Science, Xi'an Jiaotong-Liverpool University, Suzhou 215123, Jiangsu, China\\
\texttt{24107880127@nustti.edu.cn, zhipengye@nustti.edu.cn, jf@nustti.edu.cn,}\\
\texttt{qiufeng.wang@xjtlu.edu.cn}\\
$^{*}$Equal contribution. $^{\dagger}$Corresponding author: Zhipeng Ye.}

\begin{document}
\maketitle

\begin{abstract}
A frozen CLIP region explanation can be spatially plausible while a non-target
semantic foil responds more strongly than the target. We audit this failure for
the \emph{actual} target-only CCI selector using matched prompt-mean and aggregate
worst-foil margins. On COCO/OpenAI B/16, the normalized aggregate margin is
$-.342$, 74.0\% of selected regions are negative, and hardest-foil identity has
59.8\% mean pairwise agreement across held-out prompt families. A selection-stage
target-tolerant reranking at $\epsilon=.02$ switches only 2.0--2.6\% of images but
improves normalized aggregate margin by $.0150/.0108/.0118/.0087$ across COCO/VOC
and B/16/B/32. All 10,000-resample paired confidence intervals are above zero,
whereas held-out target-response intervals include zero. Gains persist with
category-disjoint and annotation-absent foils; random and target-only feasible
reranking do not explain them. Thus worst-foil debt is reproducible, while
repairability is sparse and local within frozen CCI candidate sets, without
implying a universally superior rule.
\end{abstract}
\begin{keywords}
CLIP, explainability, interventions, model auditing, worst-foil specificity
\end{keywords}

\section{Introduction}
Frozen vision-language models such as CLIP~\cite{radford2021clip} are often interpreted by ranking image regions
according to the change in image--text similarity after an intervention.
CCI~\cite{agarwal2025cci}, for example, clusters CLIP patch tokens into regions
and ranks the region that maximizes a target-score drop. A large target drop says
that a region matters for the target; it does not establish that the effect is
\emph{specific} to the target. Removing generic object evidence can reduce a
competing semantic foil even more, while bounding-box overlap remains high.
Likewise, averaging over many easy foils can hide one high-impact confounder.

We audit the declared foil-response vector of the \emph{selected} frozen
region and report target response, worst-foil specificity, and spatial locality
separately. The contribution is the post-selection audit and controlled evidence
of local repairability, not the max/CVaR operator itself~\cite{rockafellar2000cvar}.
Relevant work includes corpus-level contrastive attribution (COCOA)~\cite{lin2023cocoa},
confusion-set saliency auditing (CASE)~\cite{williamson2025case}, CLIP attribution
methods~\cite{wang2023m2ib,li2023clipsurgery,dreyer2025cliplatent}, and recent
class-disentangled/concept-contrastive analyses~\cite{li2026cdaclip,cheng2026mgaclip,visser2026contrastiveconcept}, and structured compositional caption foils such as ARO~\cite{yuksekgonul2023aro}.

Our contributions are threefold: (i) matched prompt-mean and aggregate
worst-foil estimands for the actual CCI baseline; (ii) prompt-stability,
matched-count, category-disjoint, and annotation-absent controls; and (iii) a
selection-stage target-tolerance test against actual CCI and same-feasible-set
baselines. The effect is deliberately characterized as \emph{sparse/local
repairability}: only about 2\% of selections change, but those changes reliably
improve held-out worst-foil margins.

\section{Worst-Foil Audit and Target-Tolerant Reranking}
For checkpoint $m$, let $s_m(x,c)$ be its fixed pre-activation scaled
image--text similarity. For candidate region $R$ and concept prompt $c$, define
\begin{equation}
\dropv_R(c)=s_m(x,c)-s_m(\mathcal I(x,R),c),
\label{eq:drop}
\end{equation}
where $\mathcal I$ is a frozen internal intervention. Original target-only CCI
selects
\begin{equation}
R_{\rm CCI}=\arg\max_{R\in\regions(x)}\dropv_R^{\rm sel}(t),
\label{eq:cci}
\end{equation}
using only the selection prompt. For non-target foils $\foils_t$, a worst-foil
margin compares the target response with the strongest foil response.

\paragraph{Matched held-out estimands.}
Selection uses ``a photo of a \{category\}''. Held-out evaluation uses
``a picture of the \{category\}'', ``an image containing a \{category\}'', and
``the \{category\}''. Let $e_h(c)$ denote the stored unit-norm text embedding
for held-out prompt $h$, and
$n_c=\|H^{-1}\sum_h e_h(c)\|_2$.
We normalize every class identically as
$\widetilde d_h(c)=d_h(c)/n_c$ and~define
\begin{equation}
\begin{aligned}
P_{\max}^{\rm pm}(R,t)
&=H^{-1}\!\sum_h[\widetilde d_h(t)-\max_f\widetilde d_h(f)],\\
P_{\max}^{\rm agg}(R,t)
&=H^{-1}\!\sum_h\widetilde d_h(t)-
\max_f H^{-1}\!\sum_h\widetilde d_h(f).
\end{aligned}
\label{eq:estimands}
\end{equation}
Now the pair differs only in max/mean ordering. The same paired comparisons in
raw drop space have the same direction in all four settings. $P_{\max}$ is
ontology-relative, not human semantic ground truth or a stand-alone explanation
quality score.

\paragraph{Selection-stage target tolerance.}
For each candidate define
\begin{equation}
M^{\rm sel}(R,t)=\dropv_R^{\rm sel}(t)-
\max_{f\in\foils_t}\dropv_R^{\rm sel}(f).
\end{equation}
Given tolerance $\epsilon$, form
\begin{equation}
\regions_\epsilon=\{R:\dropv_R^{\rm sel}(t)\ge
\dropv_{R_{\rm CCI}}^{\rm sel}(t)-\epsilon\},
\end{equation}
and rerank only this feasible set:
\begin{equation}
R_{\rm WF}(\epsilon)=\arg\max_{R\in\regions_\epsilon}M^{\rm sel}(R,t).
\label{eq:repair}
\end{equation}
Because $R_{\rm CCI}\in\regions_\epsilon$, the feasible set is nonempty;
any switch loses at most $\epsilon$ of the \emph{selection-stage} target drop.
This is not a held-out target guarantee. Equation~\ref{eq:repair} reads only
selection-stage tensors, and held-out prompts are used only after the region is
fixed. We use $\epsilon=.02$ as a fixed operating point, not an optimized value.
Figure~\ref{fig:method} summarizes the separation.

\section{Experimental Protocol}
\paragraph{Models, intervention, and endpoints.}
We use $K=8$ frozen CCI regions on COCO val2014~\cite{lin2014coco}
(27,708 eligible images) and VOC2007 test~\cite{everingham2010voc} (1,943), with
OpenAI CLIP B/16 and B/32. Reranking changes only which frozen candidate is
selected. The intervention blocks selected patch key/value columns in every
visual-transformer block and head while retaining the prefix token. BBox precision
is the fraction of selected patch centers inside transformed target boxes and is
reported separately because perturbation metrics depend on the intervention
operator~\cite{gomez2022metrics}. Ties use the first exact NumPy argmax.

\paragraph{Held-out controls and uncertainty.}
Selection uses one template and evaluation uses the three held-out templates in
Eq.~\ref{eq:estimands}. Cross-foil controls use ten fixed seeds
$1701,2701,\ldots,10701$ to form disjoint category folds and pool both directions;
annotation-absent evaluation removes categories annotated as present without
claiming semantic absence. Fixed 5/19-foil subsets isolate foil-cardinality effects.
Main intervals use 10,000 paired unique-image bootstrap resamples (seed 1701).
Only four provenance-valid OpenAI settings enter the direct reranking analysis;
two historical LAION settings are excluded because their text-normalization
provenance is unresolved.

\section{Results}
\paragraph{Actual CCI exhibits worst-foil debt.}
On COCO/OpenAI B/16, original CCI has normalized
$P_{\max}^{\rm agg}=-.342$ [--.357,--.327], and 74.04\% [73.56,74.61]\% of
selected regions have negative aggregate margin. Its held-out raw target
response is $.939$ [.924,.953] and bbox precision $.679$ [.674,.683], so the
failure is not equivalent to weak target response or poor localization.

\paragraph{Sparse/local repairability is reproducible.}
Table~\ref{tab:main} gives the direct actual-CCI comparison. Normalized aggregate
and prompt-mean deltas are positive with CIs above zero in all four settings;
the corresponding prompt-mean deltas are +.0150, +.0108, +.0119, and +.00867.
Held-out target-response intervals include zero in all four. Locality is not
uniformly preserved: B/16 on COCO and VOC shows small bbox decreases, while B/32
intervals include zero.

\paragraph{The gain is not merely selector--endpoint reuse.}
Figure~\ref{fig:ident} shows that the gain persists under category-disjoint
cross-foil and annotation-absent evaluation. Pooled normalized aggregate gains
are +.01052 [+.00901,+.01198], +.00728 [+.00629,+.00846], +.00705
[+.00296,+.01140], and +.00438 [+.00131,+.00784] for COCO B/16, COCO B/32,
VOC B/16, and VOC B/32. Restricting evaluation further to annotation-absent
foils remains positive on the full eligible image population: +.00939, +.00697,
+.00809, and +.00443, with all paired CIs above zero. Thus the effect does not
require reusing the same foil categories or scoring a co-occurring annotated
object. Cross-foil, annotation-absent, and matched-count controls target ontology
reuse, annotated co-occurrence, and foil-cardinality effects, respectively.

\paragraph{Worst-foil-aware reranking, not arbitrary reranking, drives the gain.}
Table~\ref{tab:baselines} shows that uniform feasible selection and target
runner-up selection are near zero or negative, while mean-foil reranking is
weaker. The fixed Max-.1 rule nearly matches the full worst-foil margin. This is
consistent with our positioning: the contribution is the audit axis and its
controlled repairability, not novelty of a max penalty.

\paragraph{Foil count affects prevalence but not the conclusion.}
Matched 5/19-foil controls remain positive, while the negative-margin rate rises
as the foil universe expands. Thus the full-ontology failure rate partly reflects
an extreme-value effect, but the reranking gain is not created solely by foil count.

\begin{figure*}[!t]
\centering
\includegraphics[width=0.88\textwidth]{figures/figure1_method_compact.pdf}
\caption{Selection-stage target-tolerant reranking and held-out audit. The selected region is fixed before held-out evaluation.}
\label{fig:method}
\end{figure*}

\begin{table*}[t]
\centering
\caption{Actual CCI vs. WF at $\epsilon=.02$. $P=P_{\max}^{\rm agg}$; brackets are paired 95\% CIs; ``Neg.'' is CCI$\to$WF negative-margin rate.}
\label{tab:main}
\small
\setlength{\tabcolsep}{1.8pt}
\renewcommand{\arraystretch}{0.98}
\begin{tabular}{lrrrrrr}
\toprule
Setting & $P_{\rm CCI}$ & $\Delta P$ & $\Delta d_t$ & $\Delta$bbox & Sw. & Neg.\\
\midrule
COCO B/16 & -.342 & +.0150 [+.0131,+.0169] & +.00005 [-.00034,+.00044] & -.00118 [-.00199,-.00036] & 2.39\% & 74.04$\to$73.89\%\\
COCO B/32 & -.304 & +.0108 [+.00945,+.0122] & -.00010 [-.00046,+.00027] & -.00080 [-.00162,+.00002] & 2.21\% & 73.91$\to$73.75\%\\
VOC B/16 & +.168 & +.0118 [+.00677,+.0183] & -.00085 [-.00227,+.00059] & -.00340 [-.00652,-.00034] & 2.57\% & 58.16$\to$57.59\%\\
VOC B/32 & +.186 & +.00867 [+.00441,+.0143] & -.00069 [-.00184,+.00040] & +.00023 [-.00288,+.00341] & 1.96\% & 55.89$\to$55.58\%\\
\bottomrule
\end{tabular}
\end{table*}

\begin{figure}[!t]
\centering
\includegraphics[width=0.90\columnwidth]{figures/figure2_identification.pdf}
\caption{Held-out normalized aggregate gains at $\epsilon=.02$; bars are paired
95\% CIs.}
\label{fig:ident}
\end{figure}

\begin{table}[!t]
\centering
\caption{Identification controls at $\epsilon=.02$; normalized aggregate $\Delta P_{\max}$ with paired 95\% CIs.}
\label{tab:ident}
\small
\setlength{\tabcolsep}{2.0pt}
\renewcommand{\arraystretch}{0.96}
\begin{tabular}{lcc}
\toprule
Setting & Cross-foil & Absent-only\\
\midrule
COCO B/16 & \begin{tabular}{@{}c@{}}+.01052\\{[+.00901,+.01198]}\end{tabular} &
\begin{tabular}{@{}c@{}}+.00939\\{[+.00819,+.01062]}\end{tabular}\\
COCO B/32 & \begin{tabular}{@{}c@{}}+.00728\\{[+.00629,+.00846]}\end{tabular} &
\begin{tabular}{@{}c@{}}+.00697\\{[+.00594,+.00804]}\end{tabular}\\
VOC B/16 & \begin{tabular}{@{}c@{}}+.00705\\{[+.00296,+.01140]}\end{tabular} &
\begin{tabular}{@{}c@{}}+.00809\\{[+.00438,+.01274]}\end{tabular}\\
VOC B/32 & \begin{tabular}{@{}c@{}}+.00438\\{[+.00131,+.00784]}\end{tabular} &
\begin{tabular}{@{}c@{}}+.00443\\{[+.00214,+.00727]}\end{tabular}\\
\bottomrule
\end{tabular}
\end{table}

\FloatBarrier
\begin{table}[!t]
\centering
\caption{Same-feasible-set baselines at $\epsilon=.02$; normalized aggregate $\Delta P_{\max}$. Max-.1: $d_t-.1\max_f d_f$.}
\label{tab:baselines}
\small
\setlength{\tabcolsep}{1.8pt}
\renewcommand{\arraystretch}{0.96}
\begin{tabular}{lrrrrr}
\toprule
Setting & WF & Max-.1 & Mean & Unif. & 2nd $d_t$\\
\midrule
COCO B/16 & .0150 & .0147 & .0025 & .0006 & .0002\\
COCO B/32 & .0108 & .0106 & .0015 & -.0025 & -.0023\\
VOC B/16  & .0118 & .0113 & .0073 & .0010 & .0026\\
VOC B/32  & .0087 & .0083 & .0055 & -.0044 & -.0046\\
\bottomrule
\end{tabular}
\end{table}

\paragraph{Reading the controls together.}
Tables~\ref{tab:main}--\ref{tab:baselines} vary the selector, evaluation ontology,
and reranking rule, respectively; their agreement argues against selector--endpoint
reuse, foil-count artifacts, and arbitrary switching.

\section{Additional Results and Conclusion}
\paragraph{Foil-cardinality sensitivity.}
The high full-ontology failure rate is partly an extreme-value effect: adding
more candidate foils makes a strong non-target response more likely. Table~\ref{tab:foilcount}
shows this directly. At the same time, the reranking gain stays positive at both
matched counts, so the effect is not created solely by evaluating against a larger
foil set.

\begin{figure}[!t]
\centering
\includegraphics[width=\columnwidth]{figures/qualitative_cases_wide.png}
\caption{COCO/OpenAI B/16 cases at $\epsilon=.02$. Red/blue denote CCI/WF selected patches; bbox precision: (a) $.46\!\to\!1.00$, (b) $.81\!\to\!.11$, (c) $.02\!\to\!.88$.}
\label{fig:cases}
\end{figure}

\paragraph{Debt persists under conventional screens.}
On COCO B/16, 41.18\% of original selections simultaneously have bbox precision
$\ge .5$, a positive aggregate mean-class margin, and a negative worst-foil margin.

\begin{table}[!t]
\centering
\caption{COCO foil-count sensitivity. ``Neg.'' is original-CCI negative-margin rate; $\Delta P$ is normalized aggregate gain.}
\label{tab:foilcount}
\small
\setlength{\tabcolsep}{2.1pt}
\renewcommand{\arraystretch}{0.96}
\begin{tabular}{lrrrrr}
\toprule
Setting & Neg.@5 & Neg.@19 & Neg.@all & $\Delta P_5$ & $\Delta P_{19}$\\
\midrule
COCO B/16 & 45.8\% & 63.5\% & 74.0\% & +.00872 & +.01074\\
COCO B/32 & 45.0\% & 62.8\% & 73.9\% & +.00841 & +.00947\\
\bottomrule
\end{tabular}
\end{table}

\paragraph{Prompt-family stability.}
Across the three held-out phrasings, hardest-foil identity agreement averages 59.80\%,
exact three-way agreement is 46.36\%, and semantic-family agreement is 53.69\%
(Table~\ref{tab:diagnostics}); wording often preserves confounder type even when
the single hardest category changes.

\paragraph{Sparse selection changes, large conditional corrections.}
On COCO B/16, 74.04\% of original margins are negative, but only 663/27,708
(2.39\%) selections switch at $\epsilon=.02$; 92.2\% of switches improve, with
mean/median aggregate $\Delta P_{\max}=+.626/+.406$. Table~\ref{tab:diagnostics}
reports the corresponding sign transitions.

\begin{table}[!h]
\centering
\caption{COCO/B16 stability and sign transitions at $\epsilon=.02$; ``repair/intro.'' is negative$\to$non-negative / reverse.}
\label{tab:diagnostics}
\small
\setlength{\tabcolsep}{3.0pt}
\renewcommand{\arraystretch}{0.94}
\begin{tabular}{lr}
\toprule
Diagnostic & Value\\
\midrule
Hardest-foil ID, prompt pair 1--2 & 59.25\%\\
Hardest-foil ID, prompt pair 1--3 & 66.82\%\\
Hardest-foil ID, prompt pair 2--3 & 53.32\%\\
Mean pairwise ID agreement & 59.80\%\\
All-three exact ID agreement & 46.36\%\\
Semantic-family agreement & 53.69\%\\
Aggregate repair / intro. & 42 / 0\\
Prompt-event repair / intro. & 128 / 3\\
\bottomrule
\end{tabular}
\end{table}

\paragraph{Predeclared tolerance frontier.}
Table~\ref{tab:frontier} uses the same frozen candidates and selection tensors;
$\epsilon=.02$ was fixed rather than selected from it. As tolerance grows, switches
increase while $\Delta P\mid\mathrm{sw}$ stays stable. Full-sample target/bbox deltas
are $(+.00005,-.00118)$ at $.02$, $(-.00049,-.00223)$ at $.05$, and
$(-.01396,-.01627)$ at $.20$, exposing the target--specificity trade-off.

\begin{table}[!t]
\centering
\caption{Predeclared COCO/B16 frontier; $\Delta P$ is full-sample gain, while $\Delta P\mid\mathrm{sw}$ and $\Pr(+)$ condition on switches.}
\label{tab:frontier}
\small
\setlength{\tabcolsep}{2.0pt}
\renewcommand{\arraystretch}{0.95}
\begin{tabular}{crrrr}
\toprule
$\epsilon$ & Sw. & $\Delta P$ & $\Delta P\mid\mathrm{sw}$ & $\Pr(+)$\\
\midrule
.01 & 1.23\% & +.00791 & +.642 & 90.9\%\\
.02 & 2.39\% & +.01497 & +.626 & 92.2\%\\
.05 & 6.24\% & +.03906 & +.626 & 91.6\%\\
.10 & 11.68\% & +.07643 & +.654 & 92.0\%\\
.20 & 21.87\% & +.14479 & +.662 & 92.9\%\\
\bottomrule
\end{tabular}
\end{table}

\paragraph{What $\epsilon$ changes.}
Across $.01$--$.20$, $\Delta P\mid\mathrm{sw}$ remains in a narrow $.626$--$.662$ band while the switch rate rises from 1.23\% to 21.87\%. Thus tolerance mainly expands the mass of locally repairable cases; the stronger target/locality cost appears only at larger operating points.

\newpage
\paragraph{Switch-conditioned repair accounting.}
Because unchanged images contribute zero selector delta, COCO/OpenAI B/16 obeys
$0.023928\times0.625624=0.014970$: the 2.39\% switch rate times the mean
switched-image correction recovers the full-sample gain. Of 663 switches, 611
improve. Table~\ref{tab:repair_accounting} shows a large conditional specificity
correction with near-zero mean target-response change; bbox remains a separate endpoint.

\begin{table}[!h]
\centering
\caption{Cross-setting repair concentration from Table~\ref{tab:main}. Conditional deltas are full-sample deltas divided by switch rate; values are approximate because Table~\ref{tab:main} is rounded.}
\label{tab:repair_accounting}
\small
\setlength{\tabcolsep}{1.7pt}
\renewcommand{\arraystretch}{0.94}
\begin{tabular}{lrrrr}
\toprule
Setting & Sw. & $\Delta P\mid\mathrm{sw}$ & $\Delta d_t\mid\mathrm{sw}$ & $\Delta$bbox$\mid\mathrm{sw}$\\
\midrule
COCO B/16 & 2.39\% & $\approx+.628$ & $\approx+.0021$ & $\approx-.049$\\
COCO B/32 & 2.21\% & $\approx+.489$ & $\approx-.0045$ & $\approx-.036$\\
VOC B/16  & 2.57\% & $\approx+.459$ & $\approx-.0331$ & $\approx-.132$\\
VOC B/32  & 1.96\% & $\approx+.442$ & $\approx-.0352$ & $\approx+.012$\\
\bottomrule
\end{tabular}
\end{table}

\paragraph{Repair is concentrated rather than ubiquitous.}
Across settings, $\Delta P\mid\mathrm{sw}\approx .44$--$.63$; target-response changes are much smaller and locality shifts remain setting dependent. Aggregate repair/reverse is 42/0 and prompt-event 128/3; most gains therefore reduce debt without crossing zero.

\paragraph{Target response and locality are separate endpoints.}
All four held-out target-response CIs include zero; bbox effects are small and
setting dependent (Table~\ref{tab:main}). Figure~\ref{fig:cases} shows per-image
non-dominance. COCO B/32 remains seed-stable: [+.00945,+.01215] vs. [+.00943,+.01221].
\paragraph{Conclusion.}
Actual CCI can carry substantial ontology-relative worst-foil debt, and the
effect survives matched estimands, disjoint/annotation-absent foils, matched foil
counts, and 10,000-resample uncertainty checks. Arbitrary feasible reranking does
not reproduce the gain. Worst-foil debt is therefore a reproducible failure mode
with sparse local repairability within frozen CCI candidate sets, while target
response and spatial locality remain separate endpoints.

\paragraph{Acknowledgment.}
Supported by the Young Scientific and Technological Talent Support Program,
Taizhou Fengcheng Talent Plan. ChatGPT assisted language/formatting and
Fig.~\ref{fig:method} layout; it generated no experimental data/results.

\clearpage
\balance
\bibliographystyle{IEEEtran}
\bibliography{references}
\end{document}
